\section{El desafío de la generalización: de Narrow a General}

\subsection{El problema de la definición del Reward Model}

Zhu Jun (朱军) señaló el desafío central que enfrenta la generalización del RL:

\begin{warningbox}{El dilema de la definición del reward}
En la demostración de teoremas matemáticos y en la programación, el reward (función de recompensa) es inequívoco: una respuesta correcta es correcta, una incorrecta es incorrecta. Pero en dominios abiertos como la conducción autónoma, la creación artística o la generación de video, los límites entre ``bueno'' y ``malo'' son difusos, y la percepción de cada persona es distinta. Cómo definir el reward model en estos dominios es el problema técnico central de la generalización del RL.
\end{warningbox}

Además, obtener datos de supervisión de proceso (process supervision) también es difícil: se requiere etiquetar cada paso del proceso de razonamiento, lo cual exige personal especializado y datos de alto valor.

\subsection{La oportunidad para las empresas emergentes}

Yang Zhilin (杨植麟) analizó, desde la perspectiva de una empresa emergente, las oportunidades que trae el paradigma O1:

\begin{knowledgebox}{Oportunidades que trae la nueva variable tecnológica}
\begin{itemize}
\item Las empresas con cierto umbral de cómputo pueden hacer innovación fundamental a nivel de algoritmos de RL, e incluso lograr avances en el modelo base
\item Las empresas con menos cómputo pueden, mediante el post-entrenamiento, alcanzar lo mejor en un dominio específico
\item Tener una dirección definida pero un camino incierto es algo bueno para las empresas emergentes
\end{itemize}
\end{knowledgebox}

\subsection{Resumen de la sección}

El desafío central de la generalización del RL radica en la definición del reward en dominios abiertos y en la obtención de datos de supervisión de proceso. Pero el camino técnico ya está claro, y combinado con modelos base sólidos, la velocidad de generalización será mayor que la de la generación anterior de RL (la era de AlphaGo).

